% !TEX root = ../main.tex
\section{Problema 4: Acertijo Lógico}

\subsection{Descripción del Problema y Abstracción CSP}
% Contexto del acertijo con tres amigos, sus edades (24, 25, 26), 
% gustos musicales (clásica, pop, jazz), y apellidos (García, González, López).

\subsubsection{Pistas del Problema}
\begin{enumerate}
    \item Juan es más viejo que González, a quien le gusta la música clásica.
    \item El fan de la música pop, que no es García, no tiene 24.
    \item Oscar, quien no es López, tiene 25.
    \item La música favorita de Darío no es el jazz.
    \item Las edades de los 3 amigos están entre 24 y 26 años.
\end{enumerate}

\subsubsection{Variables de Decisión y Dominios}
% Explicar el enfoque de modelado de variables (por persona o dual/canalización):
% Variables para edad, apellido y tipo de música, especificando sus dominios finitos.

\subsubsection{Restricciones del Modelo}
% Formalizar cada una de las 5 pistas en lenguaje matemático/lógico.
% Incluir restricciones de biyección / alldifferent para nombres, apellidos, música y edades.

\subsection{Detalles de Implementación (\texttt{acertijo.mzn})}
% Describir aspectos de la implementación en MiniZinc sin incluir código completo.

\subsubsection{Restricciones Redundantes y Rompimiento de Simetrías}
% Analizar si se requieren o si las restricciones directas bastan para encontrar la solución única.

\subsection{Estrategias de Búsqueda y Distribución}
% Anotaciones de búsqueda exploradas para la resolución del acertijo.

\subsection{Pruebas y Resultados}
% Presentar la solución única encontrada y estadísticas de búsqueda (nodos, fallos, tiempo).

\begin{table}[htbp]
    \centering
    \begin{tabular}{lcccc}
        \toprule
        Nombre & Apellido & Edad & Género Musical Favorito \\
        \midrule
        Juan & -- & -- & -- \\
        Oscar & -- & 25 & -- \\
        Darío & -- & -- & -- \\
        \bottomrule
    \end{tabular}
    \caption{Solución identificada para el Acertijo Lógico.}
    \label{tab:acertijo_solucion}
\end{table}

\subsection{Conclusiones del Ejercicio}
% Discusión sobre el espacio de búsqueda reducido y la efectividad de la propagación.
